\documentclass[12pt]{article}
\usepackage[utf8]{inputenc}
\usepackage[T1]{fontenc}
\usepackage{hyperref}
\title{Tongwe Language Resources for K-12 Curriculum Development}
\author{Compiled by Claude Sonnet 5 (claude-sonnet-5) for Andr\'as Kornai}
\date{}
\begin{document}
\maketitle

\section{Introduction}
Tongwe (Kitongwe or Sitongwe, ISO 639-3 tny) is a Bantu language (Guthrie zone F.10) of northwestern Tanzania, on the eastern shore of Lake Tanganyika. Wikipedia lists about 40,000 native speakers, citing 1999--2001 data. It is closely related to Bende, with about 90\% lexical similarity. The search found no school materials and no dictionary.

\section{Primary Educational Resources}
None found.

\section{Linguistic and Reference Materials}
\begin{itemize}
\item Wikipedia article on Tongwe: \url{https://en.wikipedia.org/wiki/Tongwe_language}
\item Omniglot page on Tongwe (Sitongwe): \url{https://omniglot.com/writing/tongwe.htm}
\end{itemize}
No dictionary or grammar was found.

\section{Religious and Community Resources}
Joshua Project lists the language as written, with Bible translation status shown only as ``Portions''. Its audio Bible and Jesus Film entries are marked ``Not Available''.
\begin{itemize}
\item Joshua Project language page: \url{https://joshuaproject.net/languages/tny}
\item Joshua Project people group page: \url{https://joshuaproject.net/people_groups/15555}
\end{itemize}

\section{Digital and Audio Resources}
Global Recordings Network lists audio recordings in Tongwe (Words of Life and Bible-teaching programmes): \url{https://globalrecordings.net/en/language/tny}

\section{Curriculum Development Context}
Tongwe is vigorous, but little written material was found. Curriculum work would start from new documentation and a standard orthography, and would need to coordinate with Tanzania's Swahili-medium school system. Further checks of Bible Society of Tanzania and SIL records are advisable.
\end{document}